\documentclass[10pt,a4paper]{amsart}
\usepackage[margin=19mm]{geometry}
\usepackage{amsmath,amssymb,mathtools}
\usepackage[scr=rsfs,cal=euler]{mathalfa}
\usepackage{xfrac}
\usepackage[unicode,hidelinks,bookmarks=false]{hyperref}
\newcommand{\ie}{i.e.\ }
\newcommand{\cf}{cf.\ }
\newcommand{\resp}{respectively\ }
\newcommand{\Subsection}{\subsection}
\providecommand{\nobreakdash}{\nobreak\hbox{-}}
\setlength{\emergencystretch}{2em}
\multlinegap0pt
\usepackage{bm}
\usepackage{bbm}
\usepackage{dsfont}
\usepackage{xspace}
\usepackage{cancel}
\usepackage{mathtools}
\usepackage[shortlabels]{enumitem}
\usepackage{amsthm}
\makeatletter
\@ifundefined{lemma}{\newtheorem{lemma}{Lemma}}{}
\@ifundefined{theorem}{\newtheorem{theorem}{Theorem}}{}
\makeatother
\makeatletter
\expandafter\ifx\csname ctheorem\endcsname\relax
\newenvironment{ctheorem}[1]
  {\renewcommand\thecthm{#1}\cthm}
  {\endcthm}
\fi
\newcommand{\mc}[1]{\mathcal{#1}}
\makeatother
\title[A multiparameter integral inequality for the dyadic maximal ]{A multiparameter integral inequality for the dyadic maximal operator and applications}
\author{Eleftherios N. Nikolidakis}
\date{Source: arXiv:1905.08091v9 (CC BY 4.0)}
\begin{document}
\maketitle

\noindent\emph{Source and licence.} A multiparameter integral inequality for the dyadic maximal operator and applications. Version arXiv:1905.08091v9 (CC BY 4.0), distributed under the Creative Commons Attribution 4.0 International licence, as recorded in the arXiv licence metadata for this exact version and shown on the abstract page https://arxiv.org/abs/1905.08091v9. Full rights notice: licenses/arxiv-1905.08091-CC-BY-4.0.txt.\par\smallskip

\noindent\textit{Editorial scope.} Counted unit: the Lemma whose source label is \texttt{lem:3p1} in the article (source0.tex lines 247--261, source label \texttt{lem:3p1}), delivered together with the article's own complete original proof of it (source0.tex lines 271--415, one \texttt{proof} environment of 145 lines). No mathematics is written, repaired or re-derived: statement and proof are the article's own text line for line. The proof is detached from the statement in the source: the article states the result at lines 247--261 and prints its proof at lines 271--415, and the proof environment's own header names the statement, so the association is the article's own. Required context retained uncounted: thm:3p1 (lines 130--135). Every definition, display and label of those lines is retained unchanged. Omitted from this package and not redistributed: 1--129, 136--246, 262--270, 416--1127. Nothing inside a retained excerpt is omitted or altered: every retained body is delivered exactly as the article's lines are written, so no omission receipt of any kind accompanies this package. Citations: the retained excerpts carry no citation command at all, so no bibliography entry of the article is transcribed and no reference is redistributed. Reference adaptation: none was needed and none was made; every reference inside the delivered unit resolves inside it, so no unresolved reference or citation is produced. Printed slips kept literal: no printed word, symbol, hypothesis or number of the counted statement or of its proof was altered. Layout and class: the article's own class is \texttt{article} with packages including amsmath, amssymb, amsthm, amsfonts, bm, enumerate, graphicx, color, hyperref, enumitem; the wrapper is the lane's pinned \texttt{amsart} editorial wrapper at 10pt on A4, which restates the theorem-like environments the retained text needs and adds the article's own macro definitions that the retained text needs (\texttt{\textbackslash mc}), with extra wrapper packages (bm, bbm, dsfont, xspace, cancel, mathtools, enumitem). The delivered numbering is the unit's own. Assets: no retained span contains a figure environment, a graphics-inclusion command, a PDF-page inclusion, a table or a TikZ picture; no image file is copied into this package, the package declares no asset dependency and no image file is redistributed with it. Licence: version arXiv:1905.08091v9 of this article is distributed under the Creative Commons Attribution 4.0 International licence, as recorded in the arXiv licence metadata for this exact version and shown on the abstract page https://arxiv.org/abs/1905.08091v9. A full rights notice is delivered at licenses/arxiv-1905.08091-CC-BY-4.0.txt. Characters: every retained excerpt is pure ASCII, so the delivered engine, XeLaTeX with its default Latin Modern text font, reports no missing character.
\begin{theorem}\label{thm:3p1}
Let $\beta\geq\gamma\geq 0$, and $K$ an arbitrary measurable subset of $X$, with measure $k\in (0,1]$. Then for every $\phi\in L^p(X,\mu)$ such that $\int_X \phi \; d\mu = f$ and $\int_X \phi^p \; d\mu = F$ the following inequality is true
\begin{multline} \label{eq:1p9}
F\geq \left[ 1-\frac{1}{(1+\gamma)^{p-1}} \right] \int_K \phi^p \; d\mu + \frac{(p-1)\beta}{(\beta+1)^p} \int_X (M_{\mathcal{T}}\phi)^p \; d\mu \\ - \frac{(p-1)\gamma}{(\beta+1)^p} \int_K (M_{\mathcal{T}}\phi)^p \; d\mu + \frac{f^p}{(\beta+1)^{p-1}}.
\end{multline}
\end{theorem}

\begin{lemma} \label{lem:3p1}
\begin{enumerate}[i)]
\item If $I, J\in S_\phi$ then either $A(\phi,J)\cap I = \emptyset$ or $J\subseteq I$.
\item If $I\in S_\phi$, then there exists $J\in C(I)$ such that $J\notin S_\phi$.
\item For every $I\in S_\phi$ we have that
\[
I \approx \underset{\substack{S_{\phi}\ni J \subseteq I\ \,}}{\bigcup} A(\phi,J).
\]
\item For every $I\in S_\phi$ we have that
\begin{gather*}
A(\phi,I) = I\setminus \underset{\substack{J\in S_\phi : J^\star = I\ }}{\bigcup} J,\ \ \text{and thus} \\
\mu(A(\phi,I)) = \mu(I) - \sum_{\substack{J\in S_\phi: J^\star=I\ }} \mu(J).
\end{gather*}
\end{enumerate}
\end{lemma}

\begin{proof}
We begin by considering a $\mc T$-good function $\phi$, satisfying $\int_X \phi \; d\mu =f$ and $\int_X \phi^p \; d\mu =F$. Let $K$ be a measurable subset of $X$, with $\mu(K)=k\in (0,1]$ and $\beta,\gamma$ such that $\beta>\gamma>0$.

By Lemma \ref{lem:3p1} we get that $F= \int_X \phi^p \; d\mu = \sum_{\substack{I\in S_\phi}} \int_{A_I}\phi^p$, where we write $A_I$ for the set $A(\phi,I), I\in S_{\phi}$. We split the set $A_I$ in two measurable subsets $B_I,\Gamma_I$ for any $I\in S_{\phi}$, where $\mu \left( B_I \right), \mu \left( \Gamma_I \right)>0$. The choice of $B_I,\Gamma_I$ will be given in the sequel. Write $\mu (A_I)=a_I$, for $I\in S_{\phi}$. For any $I\in S_{\phi}$ we search for a constant $\tau_I>0$ for which

\begin{equation}\label{eq:3p1}
\mu (I) \, \tau_I-(\beta +1)\, \sum_{\substack{J\in S_\phi \\ J^\star=I\ }} \mu(J)- (\gamma +1)\mu(B_I)= \mu (\Gamma_I),
\end{equation}

Then (\ref{eq:3p1}) in view of Lemma \ref{lem:3p1} is equivalent to
\begin{multline} \label{eq:3p2}
\mu (I) \, \tau_I-(\beta +1)\,\left( \mu(I)-\mu(A_I) \right)- (\gamma +1)\mu(B_I)= \mu (\Gamma_I)\Leftrightarrow \\
\left[ \tau_I - (\beta +1) \right]\mu (I)+ (\beta +1)\mu(B_I)+(\beta +1)\mu(\Gamma_I)-(\gamma +1)\mu(B_I)=\mu(\Gamma_I)\Leftrightarrow \\
\left[ \tau_I - (\beta +1) \right]\mu (I)+ \beta\mu(\Gamma_I)=(\gamma-\beta)\mu(B_I),
\end{multline}
We let $\mu(\Gamma_I)=k_I\,a_I$, for some $k_I\in (0,1)$, so $\mu(B_I)=(1-k_I)a_I$. Thus (\ref{eq:3p2}) becomes
\begin{multline} \label{eq:3p3}
\left[ \tau_I - (\beta +1) \right]\mu (I)=(\gamma-\beta)(1-k_I)a_I-\beta k_I\,a_I\Leftrightarrow \\
\left[ \tau_I - (\beta +1) \right]\mu (I)= \gamma(1-k_I)a_I-\beta a_I,
\end{multline}

We now set $p_I=\frac{a_I}{\mu(I)}$, for any $I\in S_{\phi}$.
Thus (\ref{eq:3p3}) gives
\begin{multline} \label{eq:3p4}
\tau_I - (\beta +1)=\gamma(1-k_I)p_I-\beta p_I\Leftrightarrow \\
\tau_I=\left( (\beta +1)-\beta p_I \right) + (1-k_I)\gamma p_I,
\end{multline}

Note that this choice of $\tau_I, I\in S_{\phi}$, immediately gives $\tau_I>0$, since $\beta>\gamma>0$ and $0<p_I\leq 1$ for any $I\in S_{\phi}$.

We write now
\begin{multline} \label{eq:3p5}
F=\sum_{\substack{I\in S_\phi}} \int_{A_I}\phi^p\; d\mu= \sum_{\substack{I\in S_\phi}} \int_{B_I}\phi^p\; d\mu+ \sum_{\substack{I\in S_\phi}} \int_{\Gamma_I} \phi^p \; d\mu \geq \\
\geq \sum_{\substack{I\in S_\phi}} \frac{\left(\int_{B_I}\phi\; d\mu\right)^p}{\mu(B_I)^{p-1}}+\sum_{\substack{I\in S_\phi}} \frac{\left(\int_{\Gamma_I}\phi\; d\mu\right)^p}{\mu(\Gamma_I)^{p-1}},
\end{multline}
in view of H{\"o}lder's inequality. We denote the first and the second sum on the right of (\ref{eq:3p5}) by $\Sigma_1$, $\Sigma_2$ respectively.
Then by (\ref{eq:3p1}) and Lemma 3.1 iv) we obtain the following
\begin{multline} \label{eq:3p6}
\Sigma_2= \sum_{\substack{I\in S_\phi}}\frac{1}{\mu(\Gamma_I)^{p-1}}\left( \int_I \phi\; d\mu - \sum_{\substack{J\in S_\phi \\ J^\star=I\ }} \int_J \phi \; d\mu - \int_{B_I}\phi\;d\mu \right)^p=\\
=\sum_{\substack{I\in S_\phi}} \frac{\left( \mu(I)y_I-\sum_{\substack{J\in S_\phi \\ J^\star=I }} \mu(J)y_J-\int_{B_I}\phi\; d\mu \right)^p}{\left( \tau_I \mu(I)-(\beta +1)\sum_{\substack{J\in S_\phi \\ J^\star=I }}\mu(J)-(\gamma+1)\mu(B_I)\right)^{p-1}},
\end{multline}
where $y_I=Av_{I}(\phi)$, for every $I\in S_{\phi}$.
Now because of H{\"o}lder's inequality in the form
\begin{equation} \label{eq:3p7}
\frac{\left(\lambda_1+\lambda_2+...+\lambda_{\nu} \right)^p}{\left( \mu_1+\mu_2+....+\mu_{\nu} \right)^{p-1}}\leq \frac{\lambda_1^p}{\mu_1^{p-1}}+\frac{\lambda_2^p}{\mu_2^{p-1}}+...+ \frac{\lambda_{\nu}^p}{\mu_{\nu}^{p-1}},
\end{equation}
where $p>1,\mu_i>0$ and $\lambda_i\geq 0$, for $i=1,2,...,\nu$, we have in view of (\ref{eq:3p6}) that:
\begin{multline} \label{eq:3p8}
\Sigma_2\geq \sum_{\substack{I\in S_\phi}}\frac{\left( \mu(I)y_I \right)^p}{\left( \tau_I \mu(I) \right)^{p-1}}- \sum_{\substack{I\in S_\phi}}\sum_{\substack{J\in S_\phi \\ J^\star=I\ }}\frac{\left( \mu(J)y_J \right)^p}{\left( (\beta +1)\mu(J) \right)^{p-1}}-\\
-\sum_{\substack{I\in S_\phi}}\frac{1}{(\gamma+1)^{p-1}} \frac{\left( \int_{B_I}\phi\; d\mu \right)^p}{\mu(B_I)^{p-1}}.
\end{multline}
By (\ref{eq:3p8}) we obtain
\begin{multline} \label{eq:3p9}
\Sigma_1+\Sigma_2\geq \left(1-\frac{1}{(1+\gamma)^{p-1}} \right)\Sigma_1+\sum_{\substack{I\in S_\phi}}\mu(I) \frac{y_I^p}{\tau_I^{p-1}}-\sum_{\substack{I\in S_\phi \\ I\neq X}} \mu(I) \frac{y_I^p}{(\beta+1)^{p-1}}=\\
=\left(1-\frac{1}{(1+\gamma)^{p-1}} \right)\Sigma_1+\frac{y_X^p}{\tau_X^{p-1}}+\sum_{\substack{I\in S_\phi \\ I\neq X}}\mu(I) y_I^p \left( \frac{1}{\tau_I^{p-1}}-\frac{1}{(\beta+1)^{p-1}} \right)=\\
=\left(1-\frac{1}{(1+\gamma)^{p-1}} \right)\Sigma_1+\frac{f^p}{\tau_X^{p-1}}+\\
+\sum_{\substack{I\in S_\phi \\ I\neq X}} \frac{a_I}{p_I} \left( \frac{1}{\left((\beta +1-\beta p_I)+(1-k_I)\gamma p_I\right)^{p-1}} - \frac{1}{(\beta +1)^{p-1}} \right) y_I^p.
\end{multline}
Note that in (\ref{eq:3p9}) we have used the properties of the correspondence $I\longrightarrow I^{\star}$, on $S_{\phi}$.

We denote now $\Sigma_3$ the sum on the right of (\ref{eq:3p9}).
Then
\begin{equation}\label{eq:3p10}
\Sigma_3 \geq \sum_{\substack{I\in S_\phi \\ I\neq X}} \frac{1}{p_I}\left[ \frac{\beta p_I-(1-k_I)\gamma p_I}{(\beta +1)^p} (p-1) \right]a_I y_I^p,
\end{equation}
because of the inequality
\begin{equation}\label{eq:3p11}
\frac{1}{\left( (\beta +1)-s \right)^{p-1}}-\frac{1}{(\beta +1)^{p-1}}\geq\frac{(p-1)s}{(\beta +1)^p},
\end{equation}
which is true for any $\beta>0$, and $s\in[0,\beta]$, by the mean value theorem on derivatives. Note that since $\beta>\gamma>0$, we have that the quantity $s=\beta p_I-(1-k_I)\gamma p_I$ is positive and less than $\beta$ so (\ref{eq:3p11}) applies in $\Sigma_3$ , and gives (\ref{eq:3p10}). Thus
\begin{multline} \label{eq:3p12}
\Sigma_3\geq (p-1)\sum_{\substack{I\in S_\phi \\ I\neq X}} \frac{\beta -\gamma}{(\beta +1)^{p}}a_I y_I^p+ (p-1)\sum_{\substack{I\in S_\phi \\ I\neq X}}\frac{k_I \gamma}{(\beta+1)^p}a_I y_I^p=\\
=(p-1)\frac{\beta-\gamma}{(\beta +1)^p}\sum_{\substack{I\in S_\phi }}a_I y_I^p -(p-1)\frac{\beta- \gamma}{(\beta +1)^p}a_Xy_X^p+\\
+ (p-1)\frac{\gamma}{(\beta +1)^p}\sum_{\substack{I\in S_\phi }}k_I a_I y_I^p-\frac{(p-1)\gamma}{(\beta +1)^p} k_Xa_X y_X^p=\\
=(p-1)\frac{\beta-\gamma}{(\beta +1)^p}\int_X \left(M_{\mc T}\phi \right)^p\; d\mu +\\
+ (p-1)\frac{\gamma}{(\beta +1)^p}\int_{\Gamma} \left(M_{\mc T}\phi \right)^p\; d\mu-\frac{(p-1)}{(\beta +1)^p} \left( (\beta-\gamma)a_X+\gamma k_X a_X \right) f^p,
\end{multline}
where we have set $\Gamma=\bigcup_{\substack{I\in S_\phi}} \Gamma_I$.

By (\ref{eq:3p9}) and (\ref{eq:3p12}) we get
\begin{multline} \label{eq:3p13}
\Sigma_1+\Sigma_2\geq \left(1-\frac{1}{(1+\gamma)^{p-1}} \right)\Sigma_1+ (p-1)\frac{\beta-\gamma}{(\beta +1)^p}\int_X \left(M_{\mc T}\phi \right)^p\; d\mu +\\
+ (p-1)\frac{\gamma}{(\beta +1)^p}\int_{\Gamma} \left(M_{\mc T}\phi \right)^p\; d\mu + \lambda_4,
\end{multline}
where
\begin{equation}\label{eq:3p14}
\lambda_4= \frac{f^p}{\tau_X^{p-1}}-\frac{(p-1)}{(\beta +1)^p}\left( (\beta - \gamma)a_X+\gamma k_X a_X \right)f^p.
\end{equation}
By definition of $\tau_X$, (\ref{eq:3p14}) gives
\begin{multline} \label{eq:3p15}
\lambda_4= f^p \left[\frac{1}{\left((\beta+1)-\beta p_X+ (1-k_X)\gamma p_X \right)^{p-1}}-(p-1)\frac{(\beta-\gamma)a_X+\gamma a_X k_X}{(\beta +1)^p} \right]=\\
=f^p \left( \frac{1}{\left( (\beta+1)-\delta \right)^{p-1}} -(p-1) \frac{\delta}{(\beta +1)^p}\right),
\end{multline}
where $\delta=(\beta-\gamma)a_X+\gamma a_X k_X$(note that we used that $p_X=a_X$).

Now because of inequality (\ref{eq:3p11}) we have that
\[
\frac{1}{\left( (\beta+1)-s \right)^{p-1}} -(p-1) \frac{s}{(\beta +1)^p}\geq \frac{1}{(\beta +1)^{p-1}}, \forall s\in[0,\beta]
\]
and note that $\delta\in (0,\beta)$, by the definition of $\delta$.

So (\ref{eq:3p15}) gives $\lambda_4 \geq \frac{f^p}{(\beta+1)^{p-1}}$. Then, by (\ref{eq:3p13}) we have
\begin{multline} \label{eq:3p16}
\Sigma_1+\Sigma_2 \geq \left( 1-\frac{1}{(1+\gamma)^{p-1}} \right)\Sigma_1+(p-1)\frac{\beta-\gamma}{(\beta+1)^p}\int_X \left( M_{\mc T}\phi \right)^p\;d\mu +\\
+ (p-1)\frac{\gamma}{(\beta+1)^p}\int_{\Gamma} \left( M_{\mc T}\phi \right)^p\;d\mu+\frac{f^p}{(\beta+1)^{p-1}}=\\
=\left( 1-\frac{1}{(1+\gamma)^{p-1}} \right)\Sigma_1+\\
+\frac{(p-1)}{(\beta+1)^p} \left[ \beta \int_X \left( M_{\mc T}\phi \right)^p\;d\mu -\gamma \int_B \left( M_{\mc T}\phi \right)^p\;d\mu \right] + \frac{f^p}{(\beta+1)^{p-1}},
\end{multline}
where $B=\bigcup_{\substack{I\in S_\phi}}B_I=X\setminus\Gamma$.

Now (\ref{eq:3p16}) gives
\begin{multline} \label{eq:3p17}
\Sigma_2+\frac{1}{(1+\gamma)^{p-1}}\Sigma_1\geq\\
\geq \frac{f^p}{(\beta+1)^{p-1}} + \frac{(p-1)}{(\beta+1)^p} \left[ \beta \int_X \left( M_{\mc T}\phi \right)^p\;d\mu -\gamma \int_B \left( M_{\mc T}\phi \right)^p\;d\mu \right] \Rightarrow\\
\frac{1}{(1+\gamma)^{p-1}} \left( \Sigma_1 + \Sigma_2 \right)+ \left( 1-\frac{1}{(1+\gamma)^{p-1}} \right)\Sigma_2\geq\\
\geq   \frac{f^p}{(\beta+1)^{p-1}}+ \frac{(p-1)}{(\beta+1)^p} \left[ \beta \int_X \left( M_{\mc T}\phi \right)^p\;d\mu -\gamma \int_B \left( M_{\mc T}\phi \right)^p\;d\mu \right].
\end{multline}
But $F\geq\Sigma_1+\Sigma_2$, and $\Sigma_2\leq \int_\Gamma \phi^p\;d\mu$, so that we conclude from (\ref{eq:3p17}) that
\begin{multline} \label{eq:3p18}
\frac{F}{(1+\gamma)^{p-1}} +  \left( 1-\frac{1}{(1+\gamma)^{p-1}} \right) \int_\Gamma \phi^p\;d\mu \geq \frac{f^p}{(\beta+1)^{p-1}}+\\
\frac{(p-1)}{(\beta+1)^p} \left[ \beta \int_X \left( M_{\mc T}\phi \right)^p\;d\mu -\gamma \int_B \left( M_{\mc T}\phi \right)^p\;d\mu \right]
\end{multline}

Now from (\ref{eq:3p18}) we immediately get
\begin{multline} \label{eq:3p19}
\frac{1}{(1+\gamma)^{p-1}}\int_B \phi^p\; d\mu + \int_\Gamma \phi^p\;d\mu \geq\\ \geq \frac{f^p}{(\beta+1)^{p-1}}+\frac{(p-1)}{(\beta+1)^p} \left[ \beta \int_X \left( M_{\mc T}\phi \right)^p\;d\mu -\gamma \int_B \left( M_{\mc T}\phi \right)^p\;d\mu \right].
\end{multline}

But the left side of (\ref{eq:3p19}) equals $F-\left(1-\frac{1}{(1+\gamma)^{p-1}} \right) \int_B \phi^p \; d\mu$, so that (\ref{eq:3p19}) becomes
\begin{multline} \label{eq:3p20}
F\geq \left( 1-\frac{1}{(1+\gamma)^{p-1}} \right) \int_B\phi^p\;d\mu + \frac{f^p}{(\beta+1)^{p-1}} +\\
+ \frac{(p-1)\beta}{(\beta+1)^p} \int_X \left( M_{\mc T}\phi \right)^p\;d\mu - \frac{(p-1)\gamma}{(\beta+1)^p} \int_B \left( M_{\mc T}\phi \right)^p\;d\mu.
\end{multline}


Inequality (\ref{eq:3p20}) is in fact true for every choice of $B$, since every measurable subset $B$, of $X$ can be written as  $B=\bigcup_{\substack{I\in S_\phi}}B_I$, where $B_I=B\cap A_I$. Then setting $\Gamma_I=A_I\setminus B_I$ and following the above proof, we obtain the validity of (\ref{eq:3p20}). Theorem \ref{thm:3p1} is thus proved for any $\phi$ which is $\mc T$-good function (replace $B$ by $K$). Note that in the above proof we have used the fact that $\mu(B_I)>0$, for every $I\in S_{\phi}$, but this can be applied (by using the fact that $(X, \mu)$ is nonatomic) to prove (\ref{eq:3p20}) even if $\mu(B_I)=0$, for some $I\in S_{\phi}$. Now if $\phi\in L^p(X,\mu)$ is arbitrary, we consider the sequence $(\phi_m)_m$, where $\phi_m=\sum_{\substack{J\in \mc T_{(m)}}} Av_J(\phi)\, \chi_J$, and we set $\Phi_m=\sum_{\substack{J\in \mc T_{(m)}}} \max \left\lbrace Av_I(\phi): J\subseteq I \in \mc T \right\rbrace \chi_I $.

Then since $Av_J(\phi)=Av_I(\phi_m)$, for any $J\in \mc T$ for which $J\subseteq I \in \mc T_{(m)}$, we immediately see that $\Phi_m= M_{\mc T}\phi_m$.

Obviously $\int_X \phi_m \; d\mu=\int_X \phi \; d\mu=f $, and we can easily see that $F_m=\int_X \phi_m^p\;d\mu \leq \int_X \phi^p \;d\mu=F$. That is $\phi_m\in L^p (X,\mu), \forall m\in \mathbb{N}$.

Additionally $\Phi_m$ converges monotonically to $M_{\mc T}\phi$. Now $\phi_m$ is $\mc T$-good for any $m\in \mathbb{N}$, so that (\ref{eq:3p20}) is true, for $\phi_m$ , and for any $B\subseteq X$ measurable. Since $M_{\mc T}\phi_m $ increases to $ M_{\mc T}\phi$ on $X$, we get
\[ \lim_m \int_X \left(M_{\mc T}\phi_m \right)^p\;d\mu = \int_X \left( M_{\mc T}\phi\right)^p \; d\mu, \] and
\[ \lim_m \int_B \left( M_{\mc T}\phi_m \right)^p\; d\mu =\int_B \left( M_{\mc T}\phi \right)^p\;d\mu, \] while by the construction of $\phi_m$, and the fact that the tree $\mc T$ differentiates $L^1(X,\mu)$ we obtain that $\phi_m\longrightarrow \phi$, $\mu$-a.e on $X$. Now since $\phi_m\leq M_{\mc T}\phi_m\leq M_{\mc T}\phi$ and $M_{\mc T}\phi\in L^p(X,\mu)$(because $\phi \in L^p(X,\mu)$), we have, using the dominated convergence theorem that $\lim_m \int_X \phi_m^p=F$ and $\lim_m \int_B \phi_m^p=\int_B \phi^p \; d\mu$. From all these facts we deduce the validity of (\ref{eq:3p20}) for general $\phi\in L^p(X,\mu)$. For $\beta=\gamma>0$,(\ref{eq:3p20}) remains true by continuity reasons.
\end{proof}

\end{document}
